\documentclass[11pt, a4paper]{article}

\usepackage[utf8]{inputenc}
\usepackage{geometry}
\geometry{left=2.5cm, right=2.5cm, top=3cm, bottom=3cm}
\usepackage{amsmath, amssymb, amsfonts}
\usepackage{graphicx}
\usepackage{booktabs}
\usepackage{caption}
\usepackage{hyperref}
\usepackage{float}
\usepackage{enumitem}
\usepackage{setspace}

\hypersetup{
    colorlinks=true,
    linkcolor=blue,
    filecolor=magenta,      
    urlcolor=cyan,
    pdftitle={A Causal and Bayesian Multi-Factor Framework for Portfolio Optimization},
}

\title{\textbf{A Causal and Bayesian Multi-Factor Framework for Portfolio Optimization: Evidence from the Toronto Stock Exchange}}

\author{\textbf{Author Name, PMP} \\
Senior Manager, Pricing and Portfolio Optimization \\
Toronto, Ontario, Canada \\
\texttt{author.name@email.com}}
\date{\today}

\begin{document}

\maketitle

\begin{abstract}
\noindent Traditional machine learning-based multi-factor stock selection models frequently encounter severe overfitting when applied to highly noisy financial time series. Furthermore, in resource- and financial-heavy markets such as the Toronto Stock Exchange (TSX), linear predictive models often fall into "value traps" and generate excessive turnover, eroding theoretical Alpha through friction costs. In this paper, we propose an end-to-end quantitative portfolio optimization framework to address these challenges. First, we implement a sector-neutralized feature engineering pipeline combined with an exponential time-decay sample weighting mechanism to mitigate macroeconomic concept drift. Second, we propose a heterogeneous ensemble model (XGBoost, LightGBM, and MLP) utilizing Monte Carlo (MC) Dropout to quantify the epistemic uncertainty of predictions. Third, we employ Double Machine Learning (DML) to isolate the pure causal effect of earnings surprises, stripping away market confounding variables. Finally, we formulate a cost-aware Black-Litterman optimization model incorporating $L_2$ regularization and a dual-threshold rebalancing band. In a 45-month walk-forward backtest on the TSX Composite Index, the proposed framework achieves a 13.41\% annualized return with a Sharpe ratio of 0.49, successfully maintaining rigorous risk parity and strict turnover control.

\vspace{0.5cm}
\noindent \textbf{Keywords:} Quantitative Finance, Portfolio Optimization, Double Machine Learning, Epistemic Uncertainty, Black-Litterman, Rebalancing Band
\end{abstract}

\section{Introduction}
The application of machine learning in systematic equity strategies has grown exponentially. However, applying these techniques to the Toronto Stock Exchange (TSX) presents unique challenges. The index is heavily concentrated in cyclical sectors (Energy, Materials) and defensive sectors (Financials). Consequently, naive predictive algorithms tend to over-allocate capital to low-volatility, low-valuation stocks, resulting in prolonged periods of stagnation—commonly known as the value trap.

Moreover, theoretical alpha generated by high-frequency rebalancing is frequently nullified by transaction costs and bid-ask spreads in real-world trading. To bridge the gap between theoretical prediction and execution, we introduce a comprehensive pricing and portfolio optimization architecture. Our framework transitions from raw predictive modeling to sequential decision-making by integrating causal inference (Double Machine Learning) for signal extraction and a dynamically constrained Black-Litterman model for portfolio execution.

\section{Methodology}

\subsection{Sector-Neutralization and Time-Decay Weighting}
To prevent the model from indiscriminately favoring structurally low-valuation sectors, absolute fundamental metrics (e.g., Price-to-Earnings) are transformed into sector-neutral relative metrics:
\begin{equation}
PE_{rel, i, t} = \frac{PE_{i, t}}{\text{Median}(PE_{S, t})}
\end{equation}
Where $S$ denotes the Global Industry Classification Standard (GICS) sector of stock $i$. To address market regime shifts (concept drift), we apply an exponential time-decay weight to historical training samples. Assuming a half-life of $H$ months, the weight $W_t$ for a sample of age $\Delta t$ is defined as:
\begin{equation}
W_t = \left(\frac{1}{2}\right)^{\frac{\Delta t}{H}}
\end{equation}
This ensures recent market dynamics dictate gradient updates while preserving large-sample diversity from historical data.

\subsection{Heterogeneous Ensemble and Epistemic Uncertainty}
Our predictive engine ensembles gradient boosting trees (XGBoost, LightGBM) with a Multi-Layer Perceptron (MLP). Tree models are highly robust to non-linear tabular data, while the MLP is leveraged to extract epistemic uncertainty via Monte Carlo (MC) Dropout. During inference, dropout layers remain active, and $K$ stochastic forward passes are executed to compute the predictive variance:
\begin{equation}
\sigma_{epi, i}^2 = \frac{1}{K} \sum_{k=1}^{K} \left(\hat{y}_{i,k} - \frac{1}{K}\sum_{j=1}^{K}\hat{y}_{i,j}\right)^2
\end{equation}
A high $\sigma_{epi}^2$ indicates model indecision, which subsequently penalizes the view confidence matrix in the portfolio optimizer.

\subsection{Causal Alpha Extraction via Double Machine Learning (DML)}
Standard predictive models treat earnings surprises as standard features, exposing them to confounding biases (e.g., broader market rallies lifting all stocks regardless of individual earnings). We utilize Double Machine Learning (Chernozhukov et al., 2018) to isolate the true causal effect ($\theta$) of the earnings surprise treatment ($T$) on excess returns ($Y$), given the confounder matrix $X$:
\begin{align}
\tilde{Y} &= Y - \mathbb{E}[Y | X] \\
\tilde{T} &= T - \mathbb{E}[T | X] \\
\theta &= (\tilde{T}^\top \tilde{T})^{-1} \tilde{T}^\top \tilde{Y}
\end{align}
This causal coefficient dynamically scales the ensemble score, preventing the model from chasing noise during earnings seasons.

\section{Portfolio Execution and Risk Management}

\subsection{Cost-Aware Black-Litterman Optimization}
The raw ensemble scores are mapped to expected excess returns (Views) using a Grinold-inspired transformation based on individual stock volatility and cross-sectional Z-scores. To mitigate the "corner solution" problem typical in mean-variance optimization, we employ the Black-Litterman (BL) framework with an $L_2$ regularization penalty to enforce diversification. 

The optimization objective maximizes the risk-adjusted return subject to strict transaction cost ($\lambda$) penalties based on the previous period's weight vector ($w_{t-1}$):
\begin{equation}
\max_{w} \quad w^\top \Pi_{BL} - \frac{\gamma}{2} w^\top \Sigma_{LW} w - \lambda \|w - w_{t-1}\|_1 - \delta \|w\|_2^2
\end{equation}
where $\Sigma_{LW}$ is the Ledoit-Wolf shrinkage covariance matrix.

\subsection{Dual-Threshold Rebalancing Band}
To further suppress unnecessary turnover, a non-linear rebalancing band is applied before optimization. Let $R_i$ and $S_i$ be the rank and score of stock $i$. A stock currently held is retained if it satisfies either the rank buffer condition ($R_i \le R_{buffer}$) or the score tolerance condition ($S_{target} - S_i \le S_{tol}$). This guarantees that capital is only reallocated when a new asset presents a statistically significant Alpha advantage over the friction costs.

\section{Empirical Results}
We backtested the proposed framework on 136 constituent stocks of the TSX Composite Index. To prevent look-ahead bias, all fundamental data was strictly aligned using Point-in-Time (PIT) filing dates with appropriate lags. 

In a 45-month walk-forward validation (incorporating simulated bid-ask spreads and market impact models), the framework completely eliminated the zero-turnover stagnation observed in baseline models. The strategy achieved a cumulative return of 60.26\% (13.41\% annualized) compared to the benchmark's flat performance over the equivalent volatile period. The dynamic volatility scaling and sector-neutralization successfully contained the maximum drawdown to -17.10\%, demonstrating strong defensive capabilities during risk-off regimes.

\section{Conclusion}
This paper presents a robust, practitioner-oriented quantitative framework that successfully bridges the gap between predictive machine learning and institutional portfolio execution. By integrating Double Machine Learning for causal inference, MC Dropout for uncertainty quantification, and a highly constrained Black-Litterman optimizer, the system effectively navigates the unique structural challenges of the Canadian equities market. Future research will focus on extending the causal inference module to incorporate alternative high-frequency datasets and evaluating the framework's scalability in broader North American markets.

\end{document}
